\minitopic{确证、保证与知识}英文文献中 justification 与 warrant 有时承担不同角色。确证可指主体从自身视角相信得合理；Plantinga用 warrant 指把真信念提升为知识的性质\parencite{plantinga1993}。若翻译都用“确证”，会把“理性但不幸为假”与“足以构成知识”混在一起。本书在一般争论中用确证，在特定理论中说明作者技术含义。 一个主体处在完美欺骗环境，严格依证据行动。他的信念系统大量为假，却似乎无可责且内部合理。内在主义倾向保留确证，外在主义可能否认知识保证。多层评价允许说：主体具有主观或内部合理性，却缺乏连接真理的客观地位。 知识还要求反运气条件，因此即使信念得到内部和外部意义的确证，也可能因个案巧合失败。把每种失败都定义成“没有真正确证”，可以在术语上保住JTB，却会使确证变得事实性并远离日常理性评价。

\minitopic{基础信念如何获得支持}古典基础主义常把基础信念设为不可错、不可怀疑或关于当前经验。温和基础主义只要求非推论支持：经验可直接给外部世界信念初步确证，同时允许后续击败。这样无需从“我有红色感觉”推理到“有红物”，也不把经验判断设为绝对可靠。 Sellars式批评担心“给予神话”：若经验只是非概念刺激，它不能成为理由；若已经带有命题内容，又似乎需要自身确证。回应可区分经验拥有内容与主体作出判断，或主张经验的理由角色来自正常知觉能力而非另一个信念。 基础主义还面对任意停止问题。为什么经验可以停止回溯，愿望不可以？答案不能只是“因为它基础”。经验与所支持命题具有现象呈现、因果联系和可靠性等特征；不同基础理论选取不同特征。停止点需要独立解释。

\minitopic{融贯不仅是一致}一组信念不矛盾，只是最低一致性。融贯还包括解释联系、概率协调、推理整合和对异常的处理。两个系统都可内部一致，但能解释更多数据、使用更少临时假设、能预测新事实者具有更强融贯。 “孤立系统”反例设想一个完全融贯却与世界脱节的信念网。融贯主义可加入经验自发性：经验不是任意选择的信念，而对系统施加约束。也可采用外在融贯，要求系统形成方式可靠。加入这些条件后，理论与基础主义或可靠主义的差距缩小。 融贯具有整体主义代价。若一个新观察与大量核心信念冲突，我们可修改观察判断而非整个系统，这常是合理的；但过强整体防御会把反证一律解释掉。健康系统需要既稳定又可修正，能够说明哪些信念处于核心、哪些异常足以触发重构。

\minitopic{证据由什么构成}证据可以包括经验、记忆、已知命题、统计记录和他人证言。若证据只包括主体当前意识内容，大量未注意却保存在记忆中的信息被排除；若包括所有心理存储，主体完全无法取用的信息又似乎不能指导理性。 “拥有证据”可能要求可在适当提示下调取，而非当下显现。考试时一时想不起公式，是否仍拥有它？答案取决于遗忘是短暂提取失败还是能力丧失。证据拥有不是单一开关，可包含可及程度与使用成本。 事实证据观强调，假的东西不能真正支持；但主体可能合理地把一份后来发现伪造的记录当证据。可区分表面证据与事实性理由，或说主体的证据是“记录看起来如此”而非记录内容本身。选择会影响对欺骗主体的评价。

\minitopic{可靠主义的一般性问题}任何具体信念都属于无限多个过程类型。一次天气判断可描述为视觉、阅读屏幕、周二阅读蓝色屏幕、在准确应用上读数等。不同类型可靠度不同，这就是\term{一般性问题}{generality problem}。理论必须说明哪种类型认识上相关。 诉诸认知心理学中的实际机制是一条路线：过程类型由产生信念的功能系统划分，而非任意语言描述。另一条路线依据反事实稳定性：类型应在近邻情形中重复。两者仍需处理机制层级与环境条件，却比事后挑选类型更有约束。 可靠度的基准类也关键。导航在城市可靠、在新修道路地区较差；是否把环境混合平均，会改变数值。知识评价通常关注主体实际环境及正常变化，而不是宇宙所有环境。环境相对性并不使可靠性主观化，但要求明确部署范围。

\minitopic{新邪恶恶魔与新知识机器}新邪恶恶魔案例中的主体与我们内部经验完全相同，却因恶魔操纵而几乎全错。许多人认为他与我们同样有确证，这反对单纯实际可靠性。可靠主义者可区分无责合理与知识导向确证，或以正常世界、设计环境而非实际短期环境评价过程。 相反案例是一台植入装置，在主体不知情时直接产生大量真信念。若装置极可靠，外在主义可能承认知识；内在主义会问主体没有理由信任突如其来的信念，何以合理。若信念伴随与正常记忆相同的现象，直觉又可能变化。案例揭示我们是否要求主体拥有关于来源的最低视角。 两类案例形成镜像：内部良好而外部失败，外部良好而内部缺失。单一条件难以兼顾。双层理论可分别评价理性与知识，但还需说明两层怎样互动，而非简单贴两个标签。

\minitopic{高阶证据与理性降级}主体有强一阶证据支持$p$，同时得知自己处于不可靠状态，如疲劳、药物影响或存在系统计算错误。高阶证据是否要求降低对$p$的信心？总是降级可能导致自我破坏：正确证明者因无根据的专家质疑而放弃；从不降级则忽视自己的可错性。 一种平衡是比较高阶证据的独立性和具体性。可靠医生确认药物显著损害当前任务，是强击败；对方只说“人人会犯错”，信息太一般。主体若能重新检查并找到错误，高阶证据被吸收；若检查能力本身受影响，则需要外部复核。 高阶证据也适用于群体。团队发现其统计流程过去有高错误率，即使当前分析看似无误，也应降低信心并审计。认识制度需要监控自身可靠性，而不只积累目标事实。

\minitopic{综合案例：历史档案争议}历史学家依据一封署名信件相信某次会议发生。基础主义分析档案知觉与鉴定报告给予的初步理由；融贯主义考察日期、人物行程与其他记录是否整合；证据主义要求态度与全部可得材料相称；可靠主义评价鉴定和保存流程；击败理论关注后来发现的伪造动机。 若信件实际为赝品，但会议由另一未发现记录证明确实发生，历史学家的真信念是Gettier式；若信件真实但他只是因愿望相信，没有阅读鉴定，可能有命题确证而无信念确证；若他认真使用当时最佳方法但信件为高明赝品，信念可内部合理却不为知识。 案例说明“有证据”不是终点。要问证据是什么、主体是否拥有、信念是否基于它、过程是否可靠、是否有击败者，以及真值是否适当源于证据。

\begin{tcolorbox}[breakable,colback=white,colframe=blue!30!black,title={确证诊断表},fonttitle=\bfseries]
分别评价：命题是否获支持；主体是否拥有支持；信念是否基于支持；有无心理或规范击败；理由对主体是否可及；过程在哪个环境与类型下可靠；一阶与高阶证据是否冲突；这些条件产生的是理性、保证还是知识。
\end{tcolorbox}
